\input{../preamble.tex}

\SetDate[09/04/2026]

\begin{document}
\pagestyle{fancy}
\fancyhead[L]{Tle STMG}
\fancyhead[C]{\textbf{Fonction inverse 2}}
\fancyhead[R]{\today}


\exe{,difficulty=1}{
	Vers quelle valeur les expressions suivantes tendent-elles lorsque $x$ tend vers $\pinfty$ ?
	\begin{multicols}{2}
	\begin{enumerate}[itemsep=10pt]
		\item
		$f(x) = \dfrac{1}x$
		\item
		$g(x) = \dfrac{1}x + 12$
		\item
		$h(x) = \dfrac{90}x$
		\item
		$i(x) = \dfrac{-5}x$
		\item
		$j(x) = 51 - \dfrac{2}x$
		\item
		$k(x) = \dfrac{2x + 61}x$
		\item
		$l(x) = \dfrac{7-x}x$
		\item
		$m(x) = \dfrac{x^2 + 2x}{x^2}  - 7$
	\end{enumerate}
	\end{multicols}
}{exe:limitpinfty}{
	
}


\exe{}{
	Trouver le nombre $x$ vérifiant chaque équation. Les dénominateurs sont supposés non nuls.
	\begin{multicols}{2}
	\begin{enumerate}[label=\alph*), itemsep=10pt]
		\item
		$\dfrac7x = 21$
		\item
		$\dfrac4{x+3} = 2$
		\item
		$-\dfrac{1,5}x + 5 = 14$
		\item
		$\dfrac{10}{2x+5} - 57 = 13$
	\end{enumerate}
	\end{multicols}
}{exe:eq-inv2}{
	todo
}

% pas convaincu par ce truc
%\exe{}{
%	En course à pied, l'allure correspond au temps nécessaire (en minutes) pour parcourir 1 kilomètre.
%	Elle s'exprime donc en min/km.
%	
%	\begin{enumerate}
%		\item
%		Une femme a mis 18 minutes pour parcourir 3 kilomètres.
%			\begin{enumerate}
%				\item
%				Quelle est son allure en min/km ?
%				\item
%				Quelle est sa vitesse moyenne en km/min ?
%				\item
%				Comment passer de l'allure à la vitesse moyenne (km/min) ? et de la vitesse moyenne (km/min) à l'allure ?
%				\item
%				Quelle est sa vitesse moyenne en km/h ?
%			\end{enumerate}
%		\item
%		Si l'allure d'une personne est de 15 min/km, quelle est sa vitesse moyenne en km/h ?
%		\item
%		Si la vitesse d'une personne est de 8km/h, quelle est son allure en min/km ?
%	\end{enumerate}
%}{exe:allure}{
%
%}


\exe{, difficulty=1}{
	Un restaurateur propose un menu unique.
		\begin{itemize}
			\item
			Les charges fixes sont estimées pour une année à 63 000€.
			\item
			Le coût de préparation d'un repas est estimé à 9€.
%			\item
%			Le prix du menu est de 16€.
		\end{itemize}
%	On suppose dans la suite que tous les repas préparés sont vendus.
		\begin{enumerate}
			\item
			Supposons que 5 000 repas soient préparés et vendus.
			\begin{enumerate}
				\item
				Donner le coût de total annuel de ces repas.
				\item
				Quel prix de repas doit être fixé pour que le restaurant devienne rentable ?
			\end{enumerate}
			\item
			Supposons désormais que $x$ repas sont préparés et vendus.
			\begin{enumerate}
				\item
				Quel prix de repas doit être fixé pour que le restaurant soit rentable ? Une fonction de $x$ est attendue.
				Remplacer $x$ par 5 000 et comparer le résultat obtenu avec la question 1 pour valider sa réponse.
				\item 
				Supposons que le restaurant prépare et vende 20 000 repas. 
				Quel prix de repas fixer pour que le restaurant soit rentable ?
			\end{enumerate}
			\item
			Combien de repas doit l'entreprise préparer et vendre pour que le prix fixé de chaque repas puisse être égal à 9,99€ tout en étant rentable ?
			\item
			Vers quelle valeur le prix de repas tend-elle lorsque le nombre de repas préparés et vendus $x$ tend vers l'infini ?
			Interpréter ce résultat en le comparant au coût de préparation d'un repas.
			\item
			Est-il possible pour le restaurant de vendre ses repas 8,99€ tout en étant rentable ?
%			\item 
%			Sous quelles conditions le restaurant peut-il baisser le prix de chaque repas ?
		\end{enumerate}
}{exe:8}{

}


%%%%%%%%%%%%

\newpage
\fancyhead[C]{\textbf{Solutions}}
\shipoutAnswer

\end{document}



